% Bagian: first-order-logic
% Bab: lindstrom
% Subbagian: lindstrom-proof

\documentclass[../../../include/open-logic-section]{subfiles}

\begin{document}

\olfileid[id]{mod}{lin}{prf}

\olsection{Teorema Lindstr\"om}

\begin{lem}
\ollabel{lem:lindstrom}
Andaikan $!E \in L(\Lang{L})$, dengan $\Lang{L}$ berhingga, dan andaikan pula
terdapat $n \in \Nat$ sedemikian sehingga bagi sembarang dua !!{structure}
$\Struct{M}$ dan~$\Struct{N}$, jika $\Struct{M} \elemequiv[n] \Struct{N}$
dan $\Struct{M} \models_L !E$, maka $\Struct{N} \models_L !E$. Maka $!E$
ekuivalen dengan !!a{sentence} orde pertama; yaitu, terdapat !!a{sentence}
orde pertama $!D$ sedemikian sehingga $\Mod(L){!E} = \Mod(L){!D}$.
\end{lem}

\begin{proof}
Ambil $n$ sedemikian sehingga sembarang dua !!{structure} yang ekuivalen
hingga peringkat~$n$, $\Struct{M}$ dan $\Struct{N}$, memberikan nilai
kebenaran yang sama kepada~$!E$. Ingat dari
\olref[bas][pis]{prop:qr-finite}: dalam !!a{language} berhingga, hanya ada
berhingga banyak kelas ekuivalensi logis !!{sentence}s orde pertama yang
berperingkat kuantor tidak lebih dari~$n$. Sekarang, untuk setiap
!!{structure} tetap~$\Struct{M}$, pilih satu wakil $!A$ dari setiap kelas
tersebut yang benar dalam~$\Struct{M}$ dan memenuhi $\QuantRank{!A}\le n$,
lalu jadikan $!D_{\Struct{M}}$ konjungsi semua wakil itu. Dengan demikian,
$\Struct{N} \models !D_{\Struct{M}}$ jika dan hanya
jika $\Struct{N} \elemequiv[n] \Struct{M}$. Selanjutnya, tetapkan
$!D = \textstyle\bigvee \Setabs{!D_{\Struct{M}}}{\Struct{M} \models_L !E}$,
dengan hanya satu wakil dipertahankan dari setiap kelas ekuivalensi logis
yang muncul dalam himpunan pada ruas kanan; jadi disjungsi ini juga berhingga.

Kesimpulan $\Mod(L){!E} = \Mod(L){!D}$ kini langsung diperoleh. Memang, jika
$\Struct{N} \models !D$, maka untuk suatu $\Struct{M} \models_L !E$ kita
mempunyai $\Struct{N} \models !D_{\Struct{M}}$, sehingga $\Struct{N}
\models_L !E$ menurut hipotesis lemma. Sebaliknya, jika $\Struct{N}
\models_L !E$, maka $!D_{\Struct{N}}$ merupakan salah satu disjung dalam
$!D$; karena $\Struct{N} \models !D_{\Struct{N}}$, kita juga mempunyai
$\Struct{N} \models !D$.
\end{proof}

\begin{thm}[Teorema Lindstr\"om]
  \ollabel{thm:lindstrom} Andaikan $\tuple{L, \models_L}$ merupakan logika
  normal yang mempunyai Sifat Kekompakan dan Sifat L\"owenheim--Skolem. Maka
  $\tuple{L, \models_L} \leq \tuple{F, \models}$ (jadi
  $\tuple{L, \models_L}$ ekuivalen dengan logika orde pertama).
\end{thm}

\begin{proof}
Menurut \olref{lem:lindstrom}, cukup ditunjukkan bahwa bagi setiap $!E
\in L(\Lang{L})$, dengan $\Lang{L}$ berhingga, terdapat $n \in \Nat$
sedemikian sehingga untuk sembarang dua !!{structure} $\Struct{M}$
dan~$\Struct{N}$ berlaku: jika $\Struct{M} \elemequiv[n] \Struct{N}$, maka
$\Struct{M}$ dan $\Struct{N}$ memberikan nilai kebenaran yang sama
kepada~$!E$. Dengan demikian $!E$ ekuivalen dengan !!a{sentence} orde pertama,
dan dari sini $\tuple{L, \models_L} \leq \tuple{F, \models}$ diperoleh.
Karena kita bekerja dalam !!a{language} relasional murni yang berhingga,
menurut \olref[bas][pis]{thm:b-n-f} kita dapat mengganti pernyataan
$\Struct{M} \elemequiv[n] \Struct{N}$ dengan pernyataan aljabar yang
bersesuaian, yakni $I_n(\emptyset,\emptyset)$.

Diberikan $!E$, andaikan demi kontradiksi bahwa untuk setiap $n$ terdapat
!!{structure}s $\Struct{M}_n$ dan $\Struct{N}_n$ sedemikian sehingga
$I_n(\emptyset, \emptyset)$, tetapi (katakanlah) $\Struct{M}_n \models_L !E$
sedangkan $\Struct{N}_n \not\models_L !E$. Berdasarkan Sifat Isomorfisme, kita
dapat mengandaikan bahwa semua $\Struct{M}_n$ menafsirkan konstanta-konstanta bahasa
itu sebagai objek-objek yang sama. Selain itu, karena hanya ada berhingga
banyak !!{sentence} atomik dalam bahasa tersebut, kita juga dapat
mengandaikan bahwa semuanya memenuhi !!{sentence}s atomik yang sama
(jika tidak, kita dapat mengambil suatu barisan bagian dari barisan
$\Struct{M}_n$). Biarkan $\Struct{M}$ menjadi gabungan semua
$\Struct{M}_n$, yaitu !!{structure} minimal tunggal yang memuat setiap
$\Struct{M}_n$ sebagai substruktur. Seperti dalam bukti
\olref[lsp]{thm:abstract-p-isom}, biarkan $\Struct{M^*}$ menjadi ekspansi
$\Struct{M}$ dengan !!{domain} $\Domain{M} \cup \Domain{M}^{<\omega}$ dalam
!!{language} diperluas yang memuat predikat-predikat konkatenasi $P$ dan~$Q$.

Dengan cara yang sama, definisikan $\Struct{N}_n$, $\Struct{N}$, dan
$\Struct{N^*}$. Sekarang, biarkan $\Struct{U}$ menjadi !!{structure} yang
!!{domain}nya mencakup !!{domain}s dari $\Struct{M^*}$ dan $\Struct{N^*}$,
serta bilangan-bilangan asli~$\Nat$ bersama urutan alaminya~$\le$, dalam
!!{language} dengan predikat-predikat tambahan yang merepresentasikan
!!{domain}s $\Domain{M}$, $\Domain{N}$, $\Domain{M}^{<\omega}$, dan
$\Domain{N}^{<\omega}$, serta predikat-predikat yang menyandikan domain dari
$\Struct{M}_n$ dan $\Struct{N}_n$ dalam arti bahwa:
\begin{align*}
  \Domain{M_n} & = \Setabs{a \in \Domain{M}}{R(a, n)}; &
  \Domain{N_n} & = \Setabs{a \in \Domain{N}}{S(a,n)}; \\
  \Domain{M}^{<\omega}_n & = \Setabs{a \in \Domain{M}^{<\omega}}{R(a,n)}; &
  \Domain{N}^{<\omega}_n & = \Setabs{a \in \Domain{N}^{<\omega}}{S(a,n)}.
\end{align*}
!!{structure}~$\Struct{U}$ juga mempunyai relasi terner $J$ sedemikian
sehingga $J(n, \mathbf{a}, \mathbf{b})$ berlaku jika dan hanya jika
$I_n(\mathbf{a}, \mathbf{b})$.

Sekarang terdapat !!a{sentence}~$!D$ dalam !!{language}~$\Lang{L}$ yang
diperluas dengan $R$, $S$, $J$, dan seterusnya, yang menyatakan bahwa $\le$
merupakan urutan linear diskret dengan unsur pertama tetapi tanpa unsur
terakhir, bahwa $\Struct{M}_n \models_L !E$ dan $\Struct{N}_n \not\models_L !E$,
serta bahwa bagi setiap $n$ dalam urutan tersebut, $J(n, \mathbf{a},
\mathbf{b})$ berlaku jika dan hanya jika $I_n(\mathbf{a}, \mathbf{b})$.

Dengan menggunakan Sifat Kekompakan, kita dapat menemukan model
$\Struct{U^*}$ dari $!D$ yang urutannya memuat suatu unsur nonstandar~$n^*$.
Secara khusus, $\Struct{U^*}$ memuat sub!!{structure}s
$\Struct{M_{n^*}}$ dan
$\Struct{N_{n^*}}$ sedemikian sehingga $\Struct{M_{n^*}} \models_L !E$ dan
$\Struct{N_{n^*}} \not\models_L !E$. Namun, sekarang kita dapat mendefinisikan
suatu himpunan $\mathcal{I}$ yang terdiri atas pasangan tupel-$k$ dari
$\Domain{M_{n^*}}$ dan $\Domain{N_{n^*}}$ dengan menetapkan
$\tuple{\mathbf{a}, \mathbf{b}} \in \mathcal{I}$ jika dan hanya jika
$J(n^*-k, \mathbf{a}, \mathbf{b})$, dengan $k$ panjang $\mathbf{a}$ dan
$\mathbf{b}$. Karena $n^*$ nonstandar, bagi setiap $k$ standar kita mempunyai
$n^* - k > 0$, dan himpunan $\mathcal{I}$ menjadi saksi bahwa
$\Struct{M_{n^*}} \simeq_p \Struct{N_{n^*}}$. Akan tetapi, menurut
\olref[lsp]{thm:abstract-p-isom}, $\Struct{M_{n^*}}$ ekuivalen-$L$ dengan
$\Struct{N_{n^*}}$, suatu kontradiksi.
\end{proof}

\end{document}
